\subsection{Góc nhìn lớp}
\label{sec:detailed-design-class-view}
\begin{figure}[H]
\centering
\IfFileExists{18_uml_class_diagram_core_domain_overview.pdf}{
  \includegraphics[width=0.97\linewidth,height=0.84\textheight,keepaspectratio]{18_uml_class_diagram_core_domain_overview.pdf}
}{
  \IfFileExists{18_uml_class_diagram_core_domain_overview.png}{
    \includegraphics[width=0.97\linewidth,height=0.84\textheight,keepaspectratio]{18_uml_class_diagram_core_domain_overview.png}
  }{
    \fbox{\parbox{0.82\linewidth}{Thiếu file \textsf{18\_uml\_class\_diagram\_core\_domain\_overview.(pdf/png)}.}}
  }
}
\caption{Sơ đồ lớp tổng quan mô tả các service ứng dụng và các điểm mở rộng chính của backend IRMS}
\label{fig:irms-class-core}
\end{figure}

Hình~\ref{fig:irms-class-core} đặt các class entrypoint và service chính cạnh các interface dùng chung, gồm \textsf{DomainEventPublisher}, \textsf{PaymentGateway}, \textsf{BillSplitStrategy}, \textsf{ReorderRecommendationStrategy} và \textsf{NotificationChannelAdapter}. 

\begin{figure}[H]
\centering
\IfFileExists{19_uml_class_diagram_ordering_and_kitchen_cluster.pdf}{
  \includegraphics[width=0.97\linewidth,height=0.84\textheight,keepaspectratio]{19_uml_class_diagram_ordering_and_kitchen_cluster.pdf}
}{
  \IfFileExists{19_uml_class_diagram_ordering_and_kitchen_cluster.png}{
    \includegraphics[width=0.97\linewidth,height=0.84\textheight,keepaspectratio]{19_uml_class_diagram_ordering_and_kitchen_cluster.png}
  }{
    \fbox{\parbox{0.82\linewidth}{Thiếu file \textsf{19\_uml\_class\_diagram\_ordering\_and\_kitchen\_cluster.(pdf/png)}.}}
  }
}
\caption{Sơ đồ lớp mô tả quan hệ giữa cụm \textsf{ordering} và \textsf{kitchen}}
\label{fig:irms-class-ordering-kitchen}
\end{figure}

Hình~\ref{fig:irms-class-ordering-kitchen} mô tả quan hệ giữa \textsf{OrdersController}, \textsf{OrdersService}, các service con của module \texttt{ordering}, cổng \textsf{KitchenCoordinationPort}, adapter \textsf{RemoteKitchenOrderRoutingClient} và phía nhận ở module \texttt{kitchen}. Quan hệ này khớp với C\&C View: module gọi món không truy cập thẳng persistence của bếp mà đi qua client tích hợp và controller nội bộ.

\begin{figure}[H]
\centering
\IfFileExists{20_uml_class_diagram_billing_inventory_and_governance_cluster.pdf}{
  \includegraphics[width=0.97\linewidth,height=0.84\textheight,keepaspectratio]{20_uml_class_diagram_billing_inventory_and_governance_cluster.pdf}
}{
  \IfFileExists{20_uml_class_diagram_billing_inventory_and_governance_cluster.png}{
    \includegraphics[width=0.97\linewidth,height=0.84\textheight,keepaspectratio]{20_uml_class_diagram_billing_inventory_and_governance_cluster.png}
  }{
    \fbox{\parbox{0.82\linewidth}{Thiếu file \textsf{20\_uml\_class\_diagram\_billing\_inventory\_and\_governance\_cluster.(pdf/png)}.}}
  }
}
\caption{Sơ đồ lớp mô tả các điểm mở rộng của cụm \textsf{billing}, \textsf{inventory} và governance}
\label{fig:irms-class-billing-inventory}
\end{figure}

Hình~\ref{fig:irms-class-billing-inventory} gom các điểm mở rộng nhạy cảm của hệ thống. Phần billing có \textsf{PaymentGateway} và các adapter thanh toán. Phần inventory có \textsf{LowStockSeverityPolicy} và \textsf{UsageBasedReorderRecommendationStrategy}. Vùng kiểm soát có \textsf{AuthService}, \textsf{AuditService}, \textsf{AuditEventMaterializer} và \textsf{IdentityRepository}.

\subsection{Giải thích theo từng cụm lớp}\label{sec:detailed-design-clusters}

\begin{figure}[H]
\centering
\IfFileExists{26_uml_class_diagram_reservation_and_seating_module.pdf}{
  \includegraphics[width=0.97\linewidth,height=0.84\textheight,keepaspectratio]{26_uml_class_diagram_reservation_and_seating_module.pdf}
}{
  \IfFileExists{26_uml_class_diagram_reservation_and_seating_module.png}{
    \includegraphics[width=0.97\linewidth,height=0.84\textheight,keepaspectratio]{26_uml_class_diagram_reservation_and_seating_module.png}
  }{
    \fbox{\parbox{0.82\linewidth}{Thiếu file \textsf{26\_uml\_class\_diagram\_reservation\_and\_seating\_module.(pdf/png)}.}}
  }
}
\caption{Sơ đồ lớp chi tiết cho module \textsf{reservation}}
\label{fig:irms-class-reservation-module}
\end{figure}

Hình~\ref{fig:irms-class-reservation-module} thể hiện các controller \textsf{ReservationController}, \textsf{TableController}, \textsf{WaitlistController}; service \textsf{ReservationService}, \textsf{ReservationCheckInService}, \textsf{ReservationWaitlistService}; workflow \textsf{ReservationSeatingMediator}; và các port/adapter JDBC cho vòng đời đặt bàn và bàn ăn. Các method trong hình được lấy từ class thật, ví dụ \textsf{load}, \textsf{recommend}, \textsf{createReservation} và \textsf{checkIn}.

\begin{figure}[H]
\centering
\IfFileExists{27_uml_class_diagram_ordering_and_menu_module.pdf}{
  \includegraphics[width=0.97\linewidth,height=0.84\textheight,keepaspectratio]{27_uml_class_diagram_ordering_and_menu_module.pdf}
}{
  \IfFileExists{27_uml_class_diagram_ordering_and_menu_module.png}{
    \includegraphics[width=0.97\linewidth,height=0.84\textheight,keepaspectratio]{27_uml_class_diagram_ordering_and_menu_module.png}
  }{
    \fbox{\parbox{0.82\linewidth}{Thiếu file \textsf{27\_uml\_class\_diagram\_ordering\_and\_menu\_module.(pdf/png)}.}}
  }
}
\caption{Sơ đồ lớp chi tiết cho module \textsf{ordering} và phần \textsf{menu}}
\label{fig:irms-class-ordering-module}
\end{figure}

Hình~\ref{fig:irms-class-ordering-module} làm rõ vai trò của \textsf{MenuItemController}, \textsf{OrdersController}, \textsf{MenuService}, \textsf{OrdersService}, \textsf{OrderCreationService}, \textsf{OrderMenuSelectionService}, các policy đơn hàng và hai phạm vi triển khaisitory port \textsf{OrderRepository}, \textsf{OrderQueryRepository}. Các adapter \textsf{JdbcOrderRepository} và \textsf{JdbcOrderQueryRepository} được đặt ở phía hiện thực port, đúng với nguyên tắc tách application service khỏi truy cập JDBC.

\begin{figure}[H]
\centering
\IfFileExists{28_uml_class_diagram_kitchen_coordination_module.pdf}{
  \includegraphics[width=0.97\linewidth,height=0.84\textheight,keepaspectratio]{28_uml_class_diagram_kitchen_coordination_module.pdf}
}{
  \IfFileExists{28_uml_class_diagram_kitchen_coordination_module.png}{
    \includegraphics[width=0.97\linewidth,height=0.84\textheight,keepaspectratio]{28_uml_class_diagram_kitchen_coordination_module.png}
  }{
    \fbox{\parbox{0.82\linewidth}{Thiếu file \textsf{28\_uml\_class\_diagram\_kitchen\_coordination\_module.(pdf/png)}.}}
  }
}
\caption{Sơ đồ lớp chi tiết cho module \textsf{kitchen}}
\label{fig:irms-class-kitchen-module}
\end{figure}

Hình~\ref{fig:irms-class-kitchen-module} mô tả các lớp điều phối bếp thật như \textsf{KitchenService}, \textsf{KitchenOrderRoutingService}, \textsf{KitchenTicketItemWorkflowService}, \textsf{KitchenWorkflowMediator}, các policy trạng thái/ưu tiên và các port \textsf{KitchenOrderRoutingRepository}, \textsf{KitchenTicketItemCommandPort}, \textsf{OrderingStatusSyncPort}. Adapter \textsf{RemoteOrderStateUpdateClient} thể hiện phụ thuộc quay về trạng thái gọi món thông qua interface thay vì truy cập trực tiếp.

\begin{figure}[H]
\centering
\IfFileExists{29_uml_class_diagram_billing_and_settlement_module.pdf}{
  \includegraphics[width=0.97\linewidth,height=0.84\textheight,keepaspectratio]{29_uml_class_diagram_billing_and_settlement_module.pdf}
}{
  \IfFileExists{29_uml_class_diagram_billing_and_settlement_module.png}{
    \includegraphics[width=0.97\linewidth,height=0.84\textheight,keepaspectratio]{29_uml_class_diagram_billing_and_settlement_module.png}
  }{
    \fbox{\parbox{0.82\linewidth}{Thiếu file \textsf{29\_uml\_class\_diagram\_billing\_and\_settlement\_module.(pdf/png)}.}}
  }
}
\caption{Sơ đồ lớp chi tiết cho module \textsf{billing}}
\label{fig:irms-class-billing-module}
\end{figure}

Hình~\ref{fig:irms-class-billing-module} thể hiện \textsf{BillingService} như facade ứng dụng. Các service con phụ trách lập hóa đơn, tách hóa đơn, thanh toán, biên nhận và hoàn tiền. Các điểm mở rộng chính là \textsf{BillSplitStrategy}, \textsf{PaymentGateway} và \textsf{ReceiptDeliveryAdapter}. Các strategy \textsf{EqualSplitStrategy}, \textsf{AmountSplitStrategy} và \textsf{SeatSplitStrategy} được biểu diễn bằng quan hệ hiện thực interface được dùng trong thiết kế.

\begin{figure}[H]
\centering
\IfFileExists{30_uml_class_diagram_inventory_and_notification_module.pdf}{
  \includegraphics[width=0.97\linewidth,height=0.84\textheight,keepaspectratio]{30_uml_class_diagram_inventory_and_notification_module.pdf}
}{
  \IfFileExists{30_uml_class_diagram_inventory_and_notification_module.png}{
    \includegraphics[width=0.97\linewidth,height=0.84\textheight,keepaspectratio]{30_uml_class_diagram_inventory_and_notification_module.png}
  }{
    \fbox{\parbox{0.82\linewidth}{Thiếu file \textsf{30\_uml\_class\_diagram\_inventory\_and\_notification\_module.(pdf/png)}.}}
  }
}
\caption{Sơ đồ lớp chi tiết cho module \textsf{inventory} và \textsf{notification}}
\label{fig:irms-class-inventory-module}
\end{figure}

Hình~\ref{fig:irms-class-inventory-module} nối \textsf{InventoryService}, \textsf{InventoryIngredientService}, \textsf{InventoryStockService} với workflow cảnh báo tồn kho và nhóm thông báo. \textsf{NotificationRequestedConsumer}, \textsf{NotificationRequestMaterializer} và các adapter kênh gửi được giữ vì chúng là class chính trong thiết kế và cũng xuất hiện ở C\&C View xử lý nền.

\begin{figure}[H]
\centering
\IfFileExists{31_uml_class_diagram_identity_audit_and_reporting_module.png}{
  \includegraphics[width=0.97\linewidth,height=0.84\textheight,keepaspectratio]{31_uml_class_diagram_identity_audit_and_reporting_module.png}
}{
  \IfFileExists{31_uml_class_diagram_identity_audit_and_reporting_module.pdf}{
    \includegraphics[width=0.97\linewidth,height=0.84\textheight,keepaspectratio]{31_uml_class_diagram_identity_audit_and_reporting_module.pdf}
  }{
    \fbox{\parbox{0.82\linewidth}{Thiếu file \textsf{31\_uml\_class\_diagram\_identity\_audit\_and\_reporting\_module.(png/pdf)}.}}
  }
}
\caption{Sơ đồ lớp chi tiết cho module \textsf{identity}, kiểm toán và \textsf{reporting}}
\label{fig:irms-class-governance-module}
\end{figure}

Hình~\ref{fig:irms-class-governance-module} mô tả vùng kiểm soát bằng các controller, service, phạm vi triển khaisitory và bộ xử lý sự kiện: \textsf{AuthController}, \textsf{AuditController}, \textsf{ReportingController}, \textsf{DashboardController}, \textsf{AuthService}, \textsf{AuditService}, \textsf{AuditEventMaterializer}, \textsf{AuditRecordingRequestedConsumer}, \textsf{ReportingQueryService} và \textsf{ReportingProjectionConsumer}. Quan hệ giữa kiểm toán/báo cáo với bộ xử lý sự kiện nhất quán với Component-and-Connector View: dữ liệu quản trị được cập nhật qua xử lý nền và phạm vi triển khaisitory/query service riêng.
\FloatBarrier
\subsection{Thiết kế hành vi}

\activitygroup{Nhóm đặt bàn và bàn ăn}

\activitydiagramfigure{uc-res-01-create-reservation}{Sơ đồ hoạt động cho UC-RES-01 -- Tạo đặt bàn}{fig:activity-uc-res-01-create-reservation}

\activitydiagramfigure{uc-res-02-checkin-and-seat}{Sơ đồ hoạt động cho UC-RES-02 -- Nhận khách và sắp bàn}{fig:activity-uc-res-02-checkin-and-seat}

\paragraph{Phân tích sơ đồ UC-RES-02.}

Luồng hoạt động trong trường hợp này bám sát thực tế tại quầy lễ tân: tìm đặt bàn, kiểm tra hiệu lực, đối chiếu số khách thực tế, xác định có cần đổi bàn hoặc ghép bàn hay không, sau đó mới mở \textsf{TableSession}. Điểm mấu chốt là việc \textbf{mở phiên phục vụ bàn} luôn diễn ra sau bước kiểm tra bàn khả dụng, nhờ đó báo cáo giữ được bất biến quan trọng rằng một bàn không có hai phiên phục vụ hoạt động cùng lúc.

Ngoài ra, sơ đồ cũng làm rõ nhánh ngoại lệ khi đặt bàn đã quá giờ giữ chỗ. Trong tình huống này, hệ thống không tự động bỏ qua mà chuyển sang bước phê duyệt của quản lý. 
\activitydiagramfigure{uc-res-03-manage-waitlist}{Sơ đồ hoạt động cho UC-RES-03 -- Quản lý danh sách chờ}{fig:activity-uc-res-03-manage-waitlist}

\activitydiagramfigure{uc-res-04-update-table-status-and-release}{Sơ đồ hoạt động cho UC-RES-04 -- Cập nhật trạng thái bàn và giải phóng bàn}{fig:activity-uc-res-04-update-table-status-and-release}

\paragraph{Phân tích sơ đồ UC-RES-04.}
\begin{itemize}
\item \textbf{Ràng buộc trạng thái \textsf{Available}:} Trạng thái bàn không cập nhật tự động sau thanh toán mà yêu cầu validation: đóng toàn bộ phiên phục vụ (\textsf{session}), tất toán công nợ (\textsf{billing}) và xác nhận hoàn tất vệ sinh (\textsf{cleaning}).
\item \textbf{Phối hợp liên module:} Đồng bộ hóa dữ liệu giữa ba miền nghiệp vụ: \textsf{TableSession}, \textsf{Billing} và \textsf{Waitlist}.
\item \textbf{Hiệu quả vận hành:} Đảm bảo tính nhất quán giữa hệ thống và thực tế, tránh lỗi gán trùng bàn (double-booking) và cung cấp dữ liệu chính xác để tính toán hiệu suất vòng quay bàn.
\end{itemize}
\activitydiagramfigure{uc-res-05-notify-customer}{Sơ đồ hoạt động cho UC-RES-05 -- Gửi thông báo cho khách hàng}{fig:activity-uc-res-05-notify-customer}

\paragraph{Phân tích sơ đồ UC-RES-05.}
\begin{itemize}
\item \textbf{Controlled Side-effect}: Thông báo là hệ quả có kiểm soát phát sinh từ các sự kiện nghiệp vụ (xác nhận đặt bàn, nhắc giờ, gọi khách từ waitlist).
\item \textbf{Quản lý luồng}: Mỗi thông báo có đầu vào, điều kiện kích hoạt và trạng thái xử lý riêng biệt.
\item \textbf{Phân tách trách nhiệm}: Cô lập logic nghiệp vụ (\textsf{Core Domain}) khỏi các chi tiết hạ tầng (\textsf{Notification Adapter}) và các cơ chế vận hành kỹ thuật (\textsf{Retry}, \textsf{Fallback}).
\end{itemize}
\activitygroup{Nhóm gọi món và điều phối bếp}

\activitydiagramfigure{uc-ord-01-create-table-order}{Sơ đồ hoạt động cho UC-ORD-01 -- Tạo đơn gọi món tại bàn}{fig:activity-uc-ord-01-create-table-order}

\paragraph{Phân tích sơ đồ UC-ORD-01.}
Luồng hoạt động gọi món bao gồm các bước: tải thực đơn, kiểm tra quyền thao tác trên phiên bàn, chọn món, kiểm tra trạng thái còn bán (quay lại chọn lại nếu hết hàng), tính tạm tính và rẽ nhánh \enquote{lưu nháp} hoặc \enquote{xác nhận đơn}. Thao tác này đóng vai trò chốt trạng thái (snapshot) về giá và tính hợp lệ của món ngay tại thời điểm xử lý.

Sơ đồ cũng vạch rõ ranh giới trách nhiệm giữa hai module:
\begin{itemize}
    \item \textbf{\textsf{Ordering}:} Giới hạn ở việc tạo và xác nhận đơn gọi món.
    \item \textbf{\textsf{Kitchen Coordination}:} Tách biệt hoàn toàn việc tạo phiếu bếp và điều phối sang trạm bếp (thuộc một sơ đồ riêng).
\end{itemize}
\activitydiagramfigure{uc-ord-02-customize-item-and-allergy-note}{Sơ đồ hoạt động cho UC-ORD-02 -- Tùy biến món và ghi chú dị ứng}{fig:activity-uc-ord-02-customize-item-and-allergy-note}

\activitydiagramfigure{uc-ord-03-send-order-to-kitchen}{Sơ đồ hoạt động cho UC-ORD-03 -- Gửi đơn gọi món sang bếp}{fig:activity-uc-ord-03-send-order-to-kitchen}

\paragraph{Phân tích sơ đồ UC-ORD-03.}
Sơ đồ mô tả ranh giới giao tiếp giữa hai module \textsf{Ordering} (gọi món) và \textsf{Kitchen Coordination} (điều phối bếp), bao gồm:

\begin{itemize}
    \item \textbf{Luồng xử lý chính:} Kiểm tra điều kiện $\rightarrow$ Phân loại món (làm ngay/làm sau) $\rightarrow$ Định tuyến theo trạm bếp (Station Routing) $\rightarrow$ Khởi tạo \textsf{KitchenTicket} $\rightarrow$ Đưa vào hàng đợi ưu tiên (Priority Queuing) $\rightarrow$ Cập nhật trạng thái đơn (\textsf{Order}).
    \item \textbf{Xử lý ngoại lệ (Exception Handling):} Nếu lỗi định tuyến trạm hoặc hệ thống màn hình bếp (KDS) offline, hệ thống đẩy request vào hàng đợi gửi lại (retry queue) và phát cảnh báo (alert) cho nhân viên. 
    \item \textbf{Đáp ứng NFR:} Cơ chế xử lý lỗi này đảm bảo yêu cầu phi chức năng về độ tin cậy (Reliability) và toàn vẹn dữ liệu (Data Integrity), tuyệt đối không làm mất các \textsf{KitchenTicket} đã được xác nhận.
\end{itemize}
\activitydiagramfigure{uc-kit-01-update-cooking-progress}{Sơ đồ hoạt động cho UC-KIT-01 -- Cập nhật tiến độ chế biến}{fig:activity-uc-kit-01-update-cooking-progress}

\paragraph{Phân tích sơ đồ UC-KIT-01.}
Luồng hoạt động bám sát quy trình xử lý thực tế tại trạm bếp:

\begin{itemize}
    \item \textbf{Luồng xử lý chính:} Nhận phiếu (\textsf{KitchenTicket}) $\rightarrow$ Kiểm tra trạm đích $\rightarrow$ Bắt đầu chế biến $\rightarrow$ Ghi nhận thời gian (timestamp) qua mỗi lần chuyển trạng thái $\rightarrow$ Đánh giá nguy cơ trễ hạn (delay risk) $\rightarrow$ Xác nhận hoàn tất (\textsf{Ready}). Việc lưu vết timestamp chi tiết là cơ sở dữ liệu bắt buộc để phục vụ module phân tích hiệu suất vận hành (operational analytics) sau này.
    
    \item \textbf{Xử lý ngoại lệ (Thiếu nguyên liệu/Sự cố):} Hệ thống không được phép kết thúc ngầm (fail silently). Request bắt buộc chuyển sang trạng thái chặn (\textsf{Blocked}), đồng thời phát cảnh báo (alert) và phản hồi trạng thái (feedback) về khu vực phục vụ. Cơ chế này đảm bảo luồng giao tiếp hai chiều giữa bếp và phục vụ luôn đồng bộ và nhất quán với kiến trúc tổng thể.
\end{itemize}
\activitydiagramfigure{uc-kit-02-serve-dishes}{Sơ đồ hoạt động cho UC-KIT-02 -- Ra món cho bàn}{fig:activity-uc-kit-02-serve-dishes}

\activitydiagramfigure{uc-kit-03-prioritize-kitchen-queue}{Sơ đồ hoạt động cho UC-KIT-03 -- Ưu tiên hàng chờ bếp và xử lý món gấp}{fig:activity-uc-kit-03-prioritize-kitchen-queue}

\paragraph{Phân tích sơ đồ UC-KIT-03.}
Luồng hoạt động định hình lại vai trò của hệ thống KDS: không chỉ hiển thị thụ động mà phải trực tiếp \textbf{điều phối thứ tự thực thi (queue prioritization)}.

\begin{itemize}
    \item \textbf{Tiêu chí phân luồng:} Trong giờ cao điểm, các \textsf{KitchenTicket} được sắp xếp tự động dựa trên: cam kết thời gian (SLA), cấu trúc món (course), phân hạng khách (VIP) và nguy cơ trễ hạn. Ca sử dụng này được tách độc lập để đáp ứng đúng yêu cầu \enquote{prioritize queues}.
    
    \item \textbf{Kiểm soát chính sách (Policy-driven):} Các can thiệp thay đổi thứ tự chế biến bắt buộc phải tuân thủ chính sách: yêu cầu phân quyền hợp lệ và ghi nhận lý do thay đổi (reason log), không cho phép tùy chỉnh tự do.
    
    \item \textbf{Khả năng truy vết (Traceability):} Cơ chế kiểm soát chặt chẽ giúp vận hành linh hoạt nhưng vẫn đảm bảo lưu vết (audit) đầy đủ để đối soát khi xảy ra khiếu nại về việc lên món sai thứ tự hoặc trễ SLA.
\end{itemize}
\activitydiagramfigure{uc-ord-04-edit-or-cancel-item}{Sơ đồ hoạt động cho UC-ORD-04 -- Sửa hoặc hủy dòng món}{fig:activity-uc-ord-04-edit-or-cancel-item}

\activitygroup{Nhóm hóa đơn, thanh toán và biên lai}

\activitydiagramfigure{uc-bil-01-create-bill}{Sơ đồ hoạt động cho UC-BIL-01 -- Tạo hóa đơn cho bàn}{fig:activity-uc-bil-01-create-bill}

\paragraph{Phân tích sơ đồ UC-BIL-01.}
\begin{itemize}
\item \textbf{Cơ chế Preview:} Bước \enquote{xem trước hóa đơn} đóng vai trò chốt chặn cuối cùng để kiểm tra tính nhất quán dữ liệu, phòng ngừa sai sót do các thay đổi tức thời (hủy món hoặc cập nhật giá).
\item \textbf{Ràng buộc nghiệp vụ:} Hệ thống thiết lập nhánh chặn phát hành khi dữ liệu không hợp lệ. Điều này khẳng định \textsf{Bill} là một thực thể nghiệp vụ (business entity) có các điều kiện phát hành và chuyển đổi trạng thái chặt chẽ, không đơn thuần là phép cộng tổng dữ liệu.
\end{itemize}
\activitydiagramfigure{uc-bil-02-split-bill-and-tip}{Sơ đồ hoạt động cho UC-BIL-02 -- Tách hóa đơn và thêm tiền boa}{fig:activity-uc-bil-02-split-bill-and-tip}

\activitydiagramfigure{uc-pay-01-process-payment}{Sơ đồ hoạt động cho UC-PAY-01 -- Xử lý thanh toán}{fig:activity-uc-pay-01-process-payment}

\paragraph{Phân tích sơ đồ UC-PAY-01.}
\begin{itemize}
    \item \textbf{Phân luồng thanh toán:}
    \begin{itemize}
        \item \textsf{PaymentGateway}: Tạo yêu cầu giao dịch $\rightarrow$ Nhận kết quả $\rightarrow$ Lưu mã tham chiếu $\rightarrow$ Xử lý các trạng thái: thành công, thất bại, hoặc chờ rà soát.
        \item \textbf{Tiền mặt}: Tính tiền thối và cập nhật số dư còn lại của hóa đơn.
    \end{itemize}
    \item \textbf{Logic tất toán:} Sau khi ghi nhận, hệ thống kiểm tra số tiền đã thanh toán:
    \begin{itemize}
        \item \textbf{Chưa đủ}: Hóa đơn giữ trạng thái mở (\textsf{OPEN}).
        \item \textbf{Đã đủ}: Đóng hóa đơn (\textsf{CLOSED}) và kích hoạt luồng phát hành biên lai.
    \end{itemize}
\end{itemize}
\activitydiagramfigure{uc-pay-02-issue-receipt}{Sơ đồ hoạt động cho UC-PAY-02 -- Phát hành biên lai}{fig:activity-uc-pay-02-issue-receipt}

\activitydiagramfigure{uc-pay-03-refund}{Sơ đồ hoạt động cho UC-PAY-03 -- Hoàn tiền}{fig:activity-uc-pay-03-refund}

\paragraph{Phân tích sơ đồ UC-PAY-03.}
Sơ đồ hoàn tiền hiện đã thể hiện đầy đủ các điểm kiểm soát bắt buộc:

\begin{itemize}
    \item Chọn giao dịch gốc
    \item Kiểm tra số tiền còn được hoàn
    \item Nhập lý do hoàn tiền
    \item So sánh với ngưỡng phê duyệt
    \item Xin xác nhận của quản lý nếu cần
    \item Xử lý hoàn tiền qua \textsf{PaymentGateway} hoặc ghi nhận hoàn tiền thủ công
    \item Cập nhật sổ kiểm toán
\end{itemize}

Điều này rất quan trọng vì hoàn tiền là thao tác nhạy cảm nhất trong vùng thanh toán.

Một chi tiết đáng chú ý là sơ đồ đã tách rõ nhánh \enquote{chờ rà soát} khi \textsf{PaymentGateway} không chấp nhận hoặc phản hồi không rõ ràng. Nhờ đó báo cáo có thể giải thích trạng thái nghiệp vụ chặt chẽ hơn, thay vì gom mọi trường hợp không thành công vào một nhãn lỗi duy nhất.
\activitygroup{Nhóm tồn kho, cảnh báo và phân tích}

\activitydiagramfigure{uc-inv-01-update-stock-consumption}{Sơ đồ hoạt động cho UC-INV-01 -- Cập nhật tiêu hao tồn kho sau phục vụ}{fig:activity-uc-inv-01-update-stock-consumption}

\paragraph{Phân tích sơ đồ UC-INV-01.}
Sơ đồ cập nhật tiêu hao tồn kho phản ánh đúng bản chất suy diễn của vùng tồn kho. Các bước xử lý bao gồm:

\begin{itemize}
    \item Nhận sự kiện từ món đã phục vụ hoặc điều chỉnh thủ công
    \item Tra công thức nguyên liệu
    \item Tính lượng tiêu hao theo số lượng món
    \item Sinh \textsf{StockTransaction} theo từng nguyên liệu
    \item Cập nhật số lượng tồn hiện có
    \item Đánh giá các điều kiện bất thường
\end{itemize}

Sơ đồ cũng làm rõ hai nhánh ngoại lệ quan trọng:

\begin{itemize}
    \item Thiếu ánh xạ công thức
    \item Âm kho bất thường
\end{itemize}

Nhờ đó phần thiết kế hành vi bám sát hơn với ca sử dụng đã đặc tả, đồng thời nối được với sơ đồ cảnh báo sắp hết hàng ngay bên dưới.
\activitydiagramfigure{uc-inv-02-low-stock-alert}{Sơ đồ hoạt động cho UC-INV-02 -- Nhận cảnh báo sắp hết hàng}{fig:activity-uc-inv-02-low-stock-alert}

\activitydiagramfigure{uc-inv-03-reorder-recommendation}{Sơ đồ hoạt động cho UC-INV-03 -- Phân tích tiêu hao và đề xuất nhập hàng}{fig:activity-uc-inv-03-reorder-recommendation}

\paragraph{Phân tích sơ đồ UC-INV-03.}

Use case này giúp phần tồn kho bám sát hơn với yêu cầu \enquote{historical usage helps predict reorder quantities and manage waste}. Nếu chỉ dừng ở cập nhật tiêu hao và cảnh báo sắp hết hàng, IRMS mới chỉ hỗ trợ phản ứng vận hành ngắn hạn, chưa đủ để ra quyết định chuẩn bị hàng hóa theo chu kỳ ngày, tuần hoặc mùa cao điểm. Do đó, sơ đồ này chuyển trọng tâm của vùng tồn kho từ \textbf{theo dõi hiện trạng} sang \textbf{hỗ trợ ra quyết định}.

Điểm đáng chú ý là sơ đồ không mô tả dự báo như một thuật toán mơ hồ mà chỉ rõ các dữ liệu đầu vào gồm tiêu hao lịch sử, cảnh báo sắp hết hàng, hao hụt, lead time và mức tồn an toàn. Nhờ đó, đề xuất nhập hàng được thể hiện là kết quả của mô hình dữ liệu và luồng xử lý đã có sẵn, không phải một chức năng bổ sung rời rạc thiếu nền dữ liệu đủ tin cậy.\activitydiagramfigure{uc-rep-01-view-revenue-and-operations-report}{Sơ đồ hoạt động cho UC-REP-01 -- Xem báo cáo doanh thu và vận hành}{fig:activity-uc-rep-01-view-revenue-and-operations-report}

\activitygroup{Nhóm quản trị cấu hình và vận hành}

\activitydiagramfigure{uc-adm-01-manage-menu-and-pricing}{Sơ đồ hoạt động cho UC-ADM-01 -- Quản lý thực đơn và giá bán}{fig:menu}
\begin{itemize}
    \item Quản lý mở màn hình cấu hình thực đơn và chọn món hoặc danh mục cần chỉnh sửa.
    
    \item Hệ thống hiển thị thông tin chi tiết của món, bao gồm giá, trạng thái và các chương trình khuyến mãi hiện tại.
    
    \item Quản lý thực hiện chỉnh sửa thông tin món và có thể áp dụng khuyến mãi nếu cần.
    
    \item Khi lưu thay đổi, quản lý gửi cấu hình món, giá, trạng thái bán, khuyến mãi hoặc công thức nguyên liệu theo form hiện tại.
    
    \item Hệ thống kiểm tra dữ liệu cấu hình:
    \begin{itemize}
        \item Nếu thiếu thông tin bắt buộc, hệ thống yêu cầu bổ sung trước khi tiếp tục.
        \item Nếu danh mục không tồn tại hoặc công thức nguyên liệu sai định dạng, hệ thống yêu cầu điều chỉnh trước khi lưu.
    \end{itemize}
    
    \item Khi dữ liệu hợp lệ, hệ thống lưu thay đổi thực đơn, công thức liên quan và ghi nhận dấu vết kiểm toán; nếu trạng thái bán thay đổi, hệ thống phát event cập nhật khả dụng.
    
    \item Các đơn hàng đang mở vẫn giữ nguyên giá đã ghi nhận trước đó để đảm bảo tính nhất quán.
\end{itemize}
\activitydiagramfigure{uc-adm-02-manage-roles-and-permissions}{Sơ đồ hoạt động cho UC-ADM-02 -- Quản lý vai trò và phân quyền}{fig:role}
\begin{itemize}
    \item Quản trị viên mở màn hình quản lý người dùng và tìm kiếm người dùng cần cập nhật quyền.
    
    \item Hệ thống hiển thị vai trò và danh sách quyền hiện tại của người dùng.
    
    \item Quản trị viên thực hiện gán hoặc điều chỉnh vai trò và quyền theo chính sách RBAC.
    
    \item Hệ thống xử lý các trường hợp mở rộng:
    \begin{itemize}
        \item Nếu không còn vai trò nào được chọn, hệ thống chặn thao tác vì mỗi nhân viên phải giữ ít nhất một vai trò.
        \item Nếu vai trò được chọn không tồn tại trong danh sách hiện hành, hệ thống chặn thao tác và yêu cầu chọn lại.
    \end{itemize}
    
    \item Hệ thống kiểm tra tính hợp lệ của chính sách phân quyền:
    \begin{itemize}
        \item Nếu phát hiện xung đột chính sách RBAC, hệ thống yêu cầu quản trị viên điều chỉnh lại quyền.
        \item Nếu thao tác dẫn đến việc gỡ bỏ quyền quản trị cuối cùng, hệ thống chặn thao tác để đảm bảo luôn tồn tại ít nhất một tài khoản quản trị.
    \end{itemize}
    
    \item Khi dữ liệu hợp lệ, hệ thống thay thế danh sách vai trò của người dùng bằng các vai trò đã chọn.
    
    \item Hệ thống yêu cầu người dùng đăng nhập lại để áp dụng quyền mới.
    
    \item Cuối cùng, hệ thống ghi nhận toàn bộ thay đổi vào nhật ký kiểm toán nhằm phục vụ truy vết và kiểm soát.
\end{itemize}
\activitydiagramfigure{uc-adm-03-review-audit-log}{Sơ đồ hoạt động cho UC-ADM-03 -- Xem nhật ký kiểm toán và truy vết thao tác nhạy cảm}{fig:log}
\begin{itemize}
    \item Người dùng có thẩm quyền mở màn hình nhật ký kiểm toán để thực hiện tra cứu.
    
    \item Hệ thống kiểm tra quyền truy cập:
    \begin{itemize}
        \item Nếu không đủ quyền, hệ thống từ chối truy cập hoặc chỉ hiển thị dữ liệu rút gọn.
        \item Nếu hợp lệ, hệ thống cho phép tiếp tục truy vấn.
    \end{itemize}
    
    \item Người dùng nhập bộ lọc tìm kiếm (thời gian, người thao tác, loại hành động, mã liên kết).
    
    \item Hệ thống truy vấn và hiển thị danh sách nhật ký kiểm toán:
    \begin{itemize}
        \item Nếu không có dữ liệu phù hợp, hệ thống thông báo và yêu cầu điều chỉnh bộ lọc.
        \item Nếu người dùng không đủ quyền, hệ thống làm mờ hoặc ẩn thông tin nhạy cảm.
    \end{itemize}
    
    \item Người dùng chọn một bản ghi để xem chi tiết, bao gồm giá trị trước và sau thay đổi, lý do và ngữ cảnh thực hiện.
    
    \item Người dùng có thể xuất báo cáo hoặc đánh dấu bản ghi cần theo dõi.
    
    \item Hệ thống ghi nhận hành vi truy cập và thao tác trên nhật ký kiểm toán nhằm phục vụ kiểm soát và truy vết sau này.
\end{itemize}
\activitydiagramfigure{uc-ops-01-manage-staff-shifts}{Sơ đồ hoạt động cho UC-OPS-01 -- Quản lý ca làm nhân viên}{fig:calamviec}
\begin{itemize}
    \item Quản lý mở màn hình lịch làm việc và chọn ngày, khu vực cùng loại ca cần phân công.
    
    \item Hệ thống hiển thị danh sách nhân viên và các ca làm việc đã có.
    
    \item Quản lý thực hiện gán nhân viên vào ca và khu vực làm việc.
    
    \item Hệ thống hỗ trợ các trường hợp mở rộng:
    \begin{itemize}
        \item Nếu thời điểm kết thúc không sau thời điểm bắt đầu, hệ thống chặn lưu ca.
        \item Nếu nhân viên đã có ca trùng khoảng thời gian, hệ thống chặn lưu ca và yêu cầu chọn lại.
    \end{itemize}
    
    \item Hệ thống kiểm tra các ràng buộc lập lịch:
    \begin{itemize}
        \item Nếu phát hiện trùng ca, hệ thống chặn lưu và yêu cầu điều chỉnh.
        \item Nếu dữ liệu hợp lệ, hệ thống lưu \textsf{ShiftAssignment}, ghi nhật ký kiểm toán và đưa thông báo nội bộ vào hàng chờ cho nhân viên.
    \end{itemize}
    
    \item Khi dữ liệu hợp lệ, quản lý lưu lịch làm việc.
    
    \item Hệ thống cập nhật \textsf{ShiftAssignment}, ghi nhật ký kiểm toán và đưa thông báo nội bộ vào hàng chờ cho nhân viên liên quan.
\end{itemize}
\pagebreak
